% =====================================================================
%  method_gmm_labeller.tex  --  PHASE V CORRECTION 4 & 5
%
%  Drop-in description of the GMM-HMM labelling procedure, including
%  the accept/reject step that Phase V Task 12B showed is load-bearing
%  for reproducibility.  Insert into Section 3 (Method).
%
%  Numbers: all STANDBY-time totals are rounded to 2 significant
%  figures and stated as conditional on the physics battery passing,
%  per Phase V Task 11C (block-bootstrap CI ±15–25%) and
%  Task 12B (labelling CV 33.7% unconditional, 3.0% conditional).
% =====================================================================

\subsection{Machine-state inference}
\label{sec:method:labelling}

\subsubsection{Preprocessing}
\label{sec:method:preprocess}

Each machine's $1$\,Hz electrical record is preprocessed in three
deterministic steps, applied identically to every machine so that
differences in downstream results are attributable to the machines,
not to the preprocessing.

\begin{enumerate}
\item \emph{Spike masking.}  For each channel, rolling-MAD outliers
  (window $11$, threshold $5\times$MAD) are flagged and excluded
  from model fitting but retained in the record; removing them would
  destroy the row-count-to-seconds correspondence that episode
  durations depend on.

\item \emph{OFF-state denoising.}  Rows with active power below
  $5$\,W are smoothed with a rolling median (window $5$) to suppress
  sensor noise in the de-energised state without touching on-state
  dynamics.

\item \emph{Feature standardisation.}  Six electrical channels are
  log$_1$p-transformed (right-skewed channels: $P, Q, S, I$) and
  z-scored jointly before any model fitting.
\end{enumerate}

\subsubsection{Gaussian hidden Markov model}
\label{sec:method:ghmm}

Emission distributions are estimated by a Gaussian mixture model
(GMM, full covariance, $k = 4$ components) fitted on a uniform
random sample of $200{,}000$ rows.  A Gaussian HMM is then
warm-started from these emissions with $\texttt{params}=\texttt{``t''}$,
so only the transition matrix is learned from $64$ contiguous
two-hour blocks spread across the record; this decouples emission
shape (a property of the marginal distribution) from transition
dynamics (a property of the sequence) and avoids the local optima
that Baum--Welch finds when initialised randomly on this
dataset~\cite{fox2011sticky}.  Viterbi decoding is applied per
contiguous segment to prevent the decoder from asserting a transition
across an acquisition gap.  A minimum-dwell constraint of $10$\,s
is applied post-decode to remove label flicker without reshaping
the partition (Section~\ref{sec:results:task4}).

\subsubsection{Physics-validation accept/reject loop}
\label{sec:method:accept_reject}

\emph{This step is load-bearing for reproducibility and was not
explicit in the Phase I--IV description; it is stated here as a
result of Phase~V Task~12B.}

After decoding, a battery of ten physics tests (T1--T10:
power-factor separation, no-load current ratio, power conservation,
dwell plausibility, cross-machine replication, and factory-schedule
agreement) is run on a strided sample of the labelled record and
aggregated into a physics score $s \in [0, 1]$.  If $s < 0.90$,
the model is re-fitted with a different seed (incremented by a
fixed prime) and the battery is re-evaluated, up to three attempts.

The motivation is empirical.  Phase~V Task~12B fitted the labeller at
ten different seeds on identical preprocessed data from pelletizer~I
and found that three seeds (30\%) converge to a wrong cluster:

\begin{itemize}
\item seed~6: the ``STANDBY'' cluster absorbed a loaded running state
  (mean power $12{,}961$\,W instead of the correct $\sim3{,}750$\,W);
\item seeds~8 and~9: the ``STANDBY'' cluster absorbed the OFF state
  (mean power $0.8$--$1.1$\,W).
\end{itemize}

The physics battery scored those three seeds at $0.60$--$0.80$ and the
seven correct seeds at $1.00$, making the accept/reject decision exact
on this record.  Conditioning on a passing seed collapses the
STANDBY-hour coefficient of variation from $33.7\%$ (all ten seeds)
to $3.0\%$ (passing seeds), and the total reported throughout this
paper---$92 \pm 3$\,h over the 155-day record---is the mean of the
seven passing seeds.

\emph{All STANDBY-time totals in this paper are therefore stated as
conditional on the physics battery passing} (i.e., $s \ge 0.90$)
\emph{and are rounded to two significant figures}, because the
block-bootstrap sampling uncertainty (Phase~V Task~11C, one-week
blocks) is $\pm 15$--$25\%$ and the conditional labelling
uncertainty is $\pm 3\%$.  Three significant figures are not
supported by the data.

The accept/reject loop is implemented in \texttt{src/labelling.py}
(\texttt{label\_record}, \texttt{physics\_score\_threshold}=0.90,
\texttt{max\_refit\_attempts}=3) and its output is stored in the
model-info dictionary (\texttt{physics\_validation\_attempts},
\texttt{physics\_validation\_seed}, \texttt{physics\_score}) for
audit and reproduction.
